\begin{papercard}{iwaya2023}
\cf{Stages} \stagetag{A} \stagetag{D} \stagetag{H}. A privacy-centred security audit of 27 top-ranked Android mental-health apps, combining technical testing with PIA requests, readability and AI-assisted policy tools. Every finding is mapped to the LINDDUN threat taxonomy. \emph{Note:} the appended copy is the arXiv preprint (v1, January 2022), so page numbers refer to it.

\cf{Research question and stance} ``What is the current privacy status of top-ranked mental health apps?'' Privacy is defined broadly to include security and data protection, with emphasis on negative impacts on data subjects \pc{iwaya2023}{3}{This study was stimulated by one research question}. Sensitivity rationale: merely being linked to a mental-health app can reveal a condition \pc{iwaya2023}{3}{the mere link of users to a given app can reveal that they might be having}.

\cf{Sampling} The google-play-scraper Node.js module searched for free English mental-health apps and returned 250, the store maximum \pc{iwaya2023}{10}{This search resulted in 250 apps as 250 is the default maximum set by}. Refinement: $\geq$100k downloads, rating $>$4 stars, Medical or Health \& Fitness category, giving 37 apps. Manual inspection with dummy accounts kept only apps that store and transmit personal data, stopping once sufficient evidence was found \pc{iwaya2023}{10}{The analysis was stopped once sufficient evidence was gathered}. Excluded: 9 apps not collecting/transmitting data, and 1 app that would expose \emph{other users'} sensitive data if analysed (an ethical exclusion). 27 remained \pc{iwaya2023}{10}{one app that would reveal sensitive informa}.

\cf{Ecosystem model} A simplified data-flow diagram of data subjects, controllers (company servers) and processors (cloud, APIs, analytics, advertisers), with threats considered per flow \pc{iwaya2023}{5}{Simplified Data Flow Diagram (DFD) for the apps}. Data flows to third parties reduce developer control and create a shared-responsibility model \pc{iwaya2023}{6}{In such cases, developers have limited control over}.

\cf{LINDDUN} Three-step method: model the system (DFD), elicit threats per DFD element using the taxonomy, manage threats \pc{iwaya2023}{6}{privacy threat analysis methodology consists of three main steps}. Categories with app examples: Linkability, Identifiability, Non-repudiation, Detectability, Disclosure of information, Unawareness, Non-compliance \pc{iwaya2023}{7}{Unawareness: the data subject is unaware of the collection}.

\cf{Procedure} (1) Static analysis with MobSF, with manual verification of dangerous permissions, insecure PRNGs/ciphers and IV use \pc{iwaya2023}{11}{Manually evaluate whether the used}. (2) Dynamic analysis in Genymotion with MobSF and Drozer (19 apps, since 8 were incompatible), inspecting traffic, Logcat/Frida logs and stored databases by visual review and keyword search \pc{iwaya2023}{12}{To identify these insecure log entries, we visually inspected all files}. (3) Server TLS checks with Qualys SSL Labs. (4) \emph{PIA requests} emailed to every developer via store contact details \pc{iwaya2023}{13}{We contacted the companies and/or developers of the studied apps based}. (5) Readability with WebFX (FK, Fog, SMOG). (6) AI-assisted policy tools: CLAUDETTE (unfair clauses) and PrivacyCheck (10 user-control and 10 GDPR questions) \pc{iwaya2023}{14}{We performed the final component of the privacy policy analysis using two}. (7) Responsible disclosure: per-app reports emailed with a 30-day response window \pc{iwaya2023}{14}{asked them to respond within 30}. (8) Each issue cross-checked threat-by-threat against the LINDDUN threat-tree catalogue (v2.0) and mapped \pc{iwaya2023}{14}{for cross-checking every single issue in the list with respect to}.

\cf{Key findings (communication-relevant)} Only 3/27 had done a PIA (2 shared it); 20 did not answer \pc{iwaya2023}{24}{companies carried out a PIA for their apps, and only two}. Policies were at college level, and only one used a layered policy \pc{iwaya2023}{25}{Interestingly, only one app provided a layered privacy policy}. CLAUDETTE found on average 2.7 unfair clauses, most often ``unilateral change'' (18) and ``contract by using'' (16) \pc{iwaya2023}{25}{The most common type of}. PrivacyCheck scored user control 59/100 and GDPR 63/100. Weakest items: sharing with law enforcement, edit/delete, breach notification, encryption at rest \pc{iwaya2023}{26}{our sample of apps scored very poorly for questions such as}. Unawareness: no policy explained targeted-advertising risks for people with mental-health problems or re-identification risks \pc{iwaya2023}{3}{apps do not explain (e.g., in the privacy policy) the risks posed by}. Onboarding consent was typically an ``I agree'' screen \pc{iwaya2023}{32}{consent process was just an initial screen presenting the pri}.

\cf{Limitations stated} Manual checks limited to flagged code. Dynamic analysis covered only 19 apps, with no payments or premium features. AI tools have limited accuracy (CLAUDETTE 78\,\%, PrivacyCheck 60\,\%), so legal analysis would be ideal \pc{iwaya2023}{33}{require legal analysis of the text made by a privacy lawyer}.

\cf{Use for this thesis} For \stagetag{H}: LINDDUN as an \emph{organising vocabulary} for capability nodes (linkability, identifiability, detectability), and the DFD actor model as a source of actor and access assumptions. For \stagetag{C}/RQ1--RQ3: the \emph{Unawareness} category directly names the gap the thesis measures (no explanation of downstream risks), and ``contract by using'' / ``unilateral change'' are reusable stance or consent codes. Methodological cautions: AI tools are used without validation on this corpus, and PIA non-response is not evidence of non-performance. Technical testing is out of scope; do not import its findings as observations.
\end{papercard}
